\section{Trabajo Futuro}
\label{sec:trabajo-futuro}

Se declara explícitamente lo que quedó fuera de alcance de este informe, en vez de omitirlo o
maquillarlo con contenido de relleno:

\begin{enumerate}[nosep]
    \item \textbf{Despliegue público en producción:} preparado completamente (Dockerfiles verificados,
    configuración de Railway, documentación de despliegue), pero la cuenta y el despliegue real requieren
    una acción del equipo humano que no se completó a la fecha de este informe.
    \item \textbf{Cobertura de pruebas automatizadas:} 1 de 12 requisitos funcionales (RF-06, generación
    de actas) no tiene prueba automatizada dedicada (RF-02 y RF-04 se cubrieron el 2026-08-29 con
    \texttt{SolicitudServiceImplTest}/\texttt{CronogramaServiceImplTest}). Cobertura general de
    24.63\,\% de líneas (JaCoCo), contra un objetivo de 60\,\%.
    \item \textbf{Aplicación real del instrumento SUS:} el cuestionario está listo; aplicarlo a un grupo
    real de al menos 5--10 usuarios (estudiantes, docentes, coordinadores) es la tarea pendiente más
    directa y de menor esfuerzo entre las listadas aquí.
    \item \textbf{Corrección de las 14 vulnerabilidades de dependencias de build restantes en el
    frontend} (5 críticas, en \texttt{sharp}/\texttt{to-ico}/\texttt{jimp}, herramientas de generación
    de íconos que no se envían al navegador) --- la vulnerabilidad crítica que sí afectaba código
    servido a producción (\texttt{@angular/core}, detectada por ZAP) se corrigió el 2026-08-29.
    \item \textbf{Revisión sistemática de literatura formal:} la Sección~\ref{sec:trabajos-relacionados}
    es una búsqueda dirigida, no una SLR completa con protocolo pre-registrado y doble revisor
    independiente, siguiendo Kitchenham y Charters \cite{kitchenham2007}.
    \item \textbf{Persistencia de archivos subidos en producción:} el filesystem del contenedor backend
    es efímero en el despliegue de Railway planeado; los PDF de anteproyectos y actas se perderían en
    cada redeploy sin un volumen persistente o almacenamiento de objetos externo.
    \item \textbf{Firma real del docente-director sobre el SRS v1.0.0:} la página de aprobación está
    lista; la firma en sí es una acción pendiente del docente, no del equipo de desarrollo.
    \item \textbf{Rediseño de la rúbrica de cobertura:} priorizar pruebas de integración reales
    (\texttt{@SpringBootTest} con base de datos real) sobre pruebas unitarias adicionales, dado que
    actualmente el proyecto no tiene ninguna prueba de integración verdadera (solo unitarias y slice
    tests \texttt{@WebMvcTest}).
\end{enumerate}

Esta lista se prioriza deliberadamente: los ítems 1--3 son alcanzables en días con el equipo humano
actuando (no requieren investigación adicional, solo ejecución de lo ya preparado); los ítems 4--8
requieren más tiempo de desarrollo o coordinación externa (el docente-director, en el caso del ítem 7).
\newpage
